% © 2026 Ing. David Jaroš — CC BY-NC-ND 4.0
%
% This work is licensed under a Creative Commons Attribution-NonCommercial-NoDerivatives
% 4.0 International License (CC BY-NC-ND 4.0).
%
% License History: Earlier drafts (up to v0.3) were released under CC BY 4.0.
% From v0.4 onward, all material is released under CC BY-NC-ND 4.0 to protect
% the integrity of the theoretical work during ongoing academic development.
%
% See LICENSE.md for full license text.
%
% Track: T1_GR — General Relativity Recovery
% Purpose: Standalone Main Theorem statement — can be \input into any document
%          or compiled independently.
% Full proof chain: GR_theorem_result.tex
% Generated: 2026-04-28

\ifdefined\INCLUDEMODE
  % Being included in another document — skip preamble
\else
  \documentclass[12pt,a4paper]{article}
  \usepackage[a4paper,margin=2.5cm]{geometry}
  \usepackage{amsmath,amssymb,amsthm}
  \usepackage{hyperref}
  \usepackage{tcolorbox}
  \tcbuselibrary{skins,breakable}
  \usepackage{enumitem}

  \newtheorem{theorem}{Theorem}
  \theoremstyle{remark}
  \newtheorem{remark}{Remark}

  \title{\textbf{Main Theorem: GR Recovery in UBT}\\[4pt]
    \large Standalone Theorem Statement\\
    \normalsize Track T1\_GR}
  \author{Ing.\ David Jaroš}
  \date{April 2026}

  \begin{document}
  \maketitle
\fi

% ============================================================
%  MAIN THEOREM — GR Recovery
% ============================================================

\begin{tcolorbox}[enhanced,colback=blue!4!white,colframe=blue!55!black,
  arc=3pt,boxrule=1.2pt,
  title={\textbf{\large Main Theorem — GR Recovery (T1\_GR)}}]

\textbf{Setup.}
Let $\mathbb{B} := \mathbb{C}\otimes_{\mathbb{R}}\mathbb{H} \cong \mathrm{Mat}(2,\mathbb{C})$
be the biquaternion algebra.
Let $\Theta : M^4 \times \mathbb{C}_\tau \to \mathbb{B}$ be a field in the
admissible class
\[
  \mathcal{A}_{\mathrm{UBT}} := \bigl\{
    \Theta \;\big|\; \{\partial_\mu\Theta\}_{\mu=0}^3
    \text{ linearly independent over } \mathbb{R} \text{ in } \mathbb{B},\;
    \Theta \neq \mathrm{const}
  \bigr\},
\]
satisfying the three UBT axioms:
\begin{itemize}[leftmargin=*,noitemsep]
  \item \textbf{AXIOM-A}: algebra $\mathbb{B} \cong \mathrm{Cl}_{1,3}(\mathbb{R})$;
  \item \textbf{AXIOM-B}: complex time $\tau = t + i\psi$; derivative $\partial_\tau$
        lies in the timelike sector $\langle\partial_\tau,\partial_\tau\rangle_\eta < 0$;
  \item \textbf{AXIOM-F}: field equation
        $\nabla^\dagger\nabla\,\Theta(q,\tau) = \kappa\,\mathcal{T}(q,\tau)$.
\end{itemize}

\medskip
\textbf{Conclusion.}  The following hold:

\begin{enumerate}[label=\textbf{(\arabic*)}]

  \item \textbf{Derived metric.}  There exists a unique real symmetric
        non-degenerate tensor on $M^4$:
        \[
          g_{\mu\nu} \;=\;
          \dfrac{\mathrm{Re}\!\bigl[\mathrm{Tr}(\partial_\mu\Theta
            \cdot\partial_\nu\Theta^\dagger)\bigr]}{\mathcal{N}},
          \qquad \mathcal{N} > 0.
        \]
        The tensor $g_{\mu\nu}$ transforms as a covariant $(0,2)$ tensor under
        diffeomorphisms.

  \item \textbf{Non-degeneracy.}  $\det(g_{\mu\nu}) \neq 0$ at every point of $M^4$.

  \item \textbf{Lorentzian signature.}  The signature of $g_{\mu\nu}$ is
        $(-,+,+,+)$ as a theorem from AXIOM-B — not a postulate.

  \item \textbf{Einstein field equations.}  The total UBT action
        \[
          S_{\mathrm{total}}[g,\Theta]
          = \frac{1}{16\pi G}\int\!\!\sqrt{-g}\,R\,\mathrm{d}^4x
          + S_\Theta[g,\Theta]
        \]
        under Hilbert variation $\delta S_{\mathrm{total}}/\delta g^{\mu\nu} = 0$ gives
        \[
          G_{\mu\nu} \;=\; 8\pi G\,T_{\mu\nu},
        \]
        where $T_{\mu\nu}$ is symmetric and satisfies $\nabla^\mu T_{\mu\nu} = 0$.

  \item \textbf{Schwarzschild solution.}  The ansatz
        $\Theta_0 = e^{i\Phi(r)}[f(r)\,\mathbf{1} + g(r)\,\boldsymbol{e}_r]$
        (unique spherically symmetric vacuum solution) reproduces the Schwarzschild
        metric in isotropic coordinates:
        \[
          g_{tt} = -\Phi(r)^2, \quad
          g_{ij} = \Psi(r)^4\,\delta_{ij}, \quad
          \Psi(r) = 1 + \tfrac{M}{2r},
        \]
        verified numerically to floating-point precision.

  \item \textbf{Linearised gravity.}  Linearisation around flat background
        reproduces the linearised Einstein equations; the odd-parity graviton
        satisfies the Regge-Wheeler equation without additional input.

\end{enumerate}

\smallskip
\textbf{Open problems} (do not affect claims (1)--(6)):
The even-parity (Zerilli) graviton equation and the off-shell $\Theta$-only
closure are open at level [L2].  See \texttt{proof\_gap\_list.md}.

\end{tcolorbox}

\bigskip
\begin{remark}[Proof status]
Claims (1)--(6) are proved at level [L1].  Full proofs are in
\texttt{GR\_theorem\_result.tex} (Section~2) with supporting files in
\texttt{canonical/gr\_closure/}.
\end{remark}

\begin{remark}[What is new]
Unlike prior biquaternion gravity approaches, UBT derives the metric (claim~(1)),
the non-degeneracy (claim~(2)), and the Lorentzian signature (claim~(3)) as
theorems.  No free parameters appear in claims (1)--(4).
\end{remark}

\ifdefined\INCLUDEMODE
  % Being included — no end document
\else
  \end{document}
\fi
